% !TEX TS-program = pdflatex
% !TEX encoding = UTF-8 Unicode

% This file is a template using the "beamer" package to create slides for a talk or presentation
% - Giving a talk on some subject.
% - The talk is between 15min and 45min long.
% - Style is ornate.

% MODIFIED by Jonathan Kew, 2008-07-06
% The header comments and encoding in this file were modified for inclusion with TeXworks.
% The content is otherwise unchanged from the original distributed with the beamer package.

\documentclass{beamer}
\setbeamercovered{invisible}



% Copyright 2004 by Till Tantau <tantau@users.sourceforge.net>.
%
% In principle, this file can be redistributed and/or modified under
% the terms of the GNU Public License, version 2.
%
% However, this file is supposed to be a template to be modified
% for your own needs. For this reason, if you use this file as a
% template and not specifically distribute it as part of a another
% package/program, I grant the extra permission to freely copy and
% modify this file as you see fit and even to delete this copyright
% notice. 


\mode<presentation>
{
  \usetheme{Warsaw}
  % or ...

  % or whatever (possibly just delete it)
}


\usepackage[english]{babel}
% or whatever

\usepackage[utf8]{inputenc}
\usepackage{qrcode}
\usepackage{hyperref}
\usepackage{ytableau}
\usepackage{longtable}
% or whatever

\usepackage{times}
\usepackage[T1]{fontenc}
\usepackage{tikz}
% Or whatever. Note that the encoding and the font should match. If T1
% does not look nice, try deleting the line with the fontenc.

%%% MATH RELATED

\title{MATH 180 Strategies of Problem Solving} 
\subtitle
{} % (optional)

\author[W.R. Casper] % (optional, use only with lots of authors)
{W.R. Casper}
% - Use the \inst{?} command only if the authors have different
%   affiliation.

\institute[California State University Fullerton] % (optional, but mostly needed)
{
  Department of Mathematics\\
  California State University Fullerton}
% - Use the \inst command only if there are several affiliations.
% - Keep it simple, no one is interested in your street address.

\subject{Talks}
% This is only inserted into the PDF information catalog. Can be left
% out. 



% If you have a file called "university-logo-filename.xxx", where xxx
% is a graphic format that can be processed by latex or pdflatex,
% resp., then you can add a logo as follows:

% \pgfdeclareimage[height=0.5cm]{university-logo}{university-logo-filename}
% \logo{\pgfuseimage{university-logo}}



% Delete this, if you do not want the table of contents to pop up at
% the beginning of each subsection:
\AtBeginSubsection[]
{
  \begin{frame}<beamer>{Outline}
    \tableofcontents[currentsection,currentsubsection]
  \end{frame}
}


% If you wish to uncover everything in a step-wise fashion, uncomment
% the following command: 

%\beamerdefaultoverlayspecification{<+->}


\begin{document}

\begin{frame}
  \titlepage
\end{frame}

\begin{frame}{Outline}
  \tableofcontents
  % You might wish to add the option [pausesections]
\end{frame}

% Since this a solution template for a generic talk, very little can
% be said about how it should be structured. However, the talk length
% of between 15min and 45min and the theme suggest that you stick to
% the following rules:  

% - Exactly two or three sections (other than the summary).
% - At *most* three subsections per section.
% - Talk about 30s to 2min per frame. So there should be between about
%   15 and 30 frames, all told.

\section{SoPS Lecture 12}

\begin{frame}{Review}
  \begin{block}{Last time}
    \begin{itemize}
\pause
      \item mathematics as a creative endeavor
      \item problem solving strategies
      \item creativity in problem solving
    \end{itemize}
  \end{block}
\end{frame}

\subsection{Starter problem}

\begin{frame}{An old problem}
  \begin{block}{Problem}
    What is the sum of the numbers
    $$1 + 2 + 3 + 4 + \dots + 99 + 100?$$
  \end{block}
\pause
  \begin{block}{Theorem}
    For any positive integer $n$,
    $$1 + 2 + 3 + 4 + \dots + (n-1) + n = \frac{n(n+1)}{2}.$$
  \end{block}
\end{frame}

\begin{frame}{2024 IMO Shortlist}

Let $\lfloor x\rfloor$ be the \textbf{floor function}, ie. what we get by rounding $x$ down to the nearest integer.  So for example

$$\lfloor 2.6\rfloor = 2,\quad \lfloor 3\rfloor  = 3,\quad \lfloor -7.1\rfloor = -8.$$

\pause
For which real numbers $x$ is the integer 

$$\lfloor x \rfloor + \lfloor 2x \rfloor + \lfloor 3x\rfloor + \dots + \lfloor nx\rfloor$$

divisible by $n$ for all positive integers $n$?
\end{frame}

\subsection{Strategies of Problem Solving}

\begin{frame}{Discussion}
  \begin{block}{Question}
    What strategies of problem solving did you use on the previous problem?
  \end{block}
\end{frame}

\begin{frame}{Putnam Problem}
  \begin{block}{Problem}
A grasshopper starts at the origin in the coordinate plane and makes a sequence of
hops. Each hop has length 5, and after each hop the grasshopper is at a point whose
coordinates are both integers; thus, there are 12 possible locations for the grasshopper
after the first hop. What is the smallest number of hops needed for the grasshopper
to reach the point (2021, 2021)
  \end{block}
\end{frame}

\subsection{Wrapping up}

\begin{frame}{Final thoughts}
  \begin{block}{Reflection}
    \begin{itemize}
\pause
      \item Look up an example of a Putnam Exam problem that looks interesting.  What is the problem statement?  What is hard about it?
\pause
      \item Write down some ideas about how you might begin to approach solving the problem you mentioned.
    \end{itemize}
  \end{block}
\begin{center}
Submit your responses by \textbf{Tonight!}
\end{center}
\end{frame}


\end{document}


